% critique.tex — f163v, vv. 80–122 : commentaire critique (articles textuels du commentaire de l'auteur, repris mot pour mot ; voir repartition.md)
\cc{80}{ἐπεὶ τρίτη 'στὶ Θεσμοφορίων ἡ μέση.}{Le manuscrit de Ravenne donne au vers 80 : « puisque c'est le troisième [jour] des Thesmophories, celui du milieu. » Problème : si les Thesmophories duraient bien trois jours, celui du milieu ne peut en aucune façon être le troisième. Cette contradiction eut la fâcheuse conséquence de déchaîner des torrents de commentaires (à commencer par celui du scholiaste) ; on a inventé, en contradiction avec tous les témoignages\nc{Outre le scholiaste, qui \italique{sait} que les Thesmophories ne duraient que trois jours – sinon il ne prendrait pas la peine de tant ergoter -- et la glose dubitative d'Hésychius (τ 1440), citons : Alciphron 2,37.2 ; 3,39 ; et la « \anglais{speculative restauration} » du papyrus n° 3540.1-2, \italique{The Oxyrhynchus Papyri}, L (1983), ed. E. W. Handley}, des Thesmophories de quatre jours à grand renfort d'érudition et de savantes notes, et ce sous le prétexte qu'il ait pu y avoir avant ces trois jours d'autres festivités ailleurs (surprenant, non ? mais il suffit de prendre un calendrier des fêtes grecques pour constater qu'avec ce genre de raisonnements, on pourrait tout aussi bien \italique{prouver} que les Thesmophories duraient 365 jours par an ; au reste, il faudrait encore montrer que le milieu d'un ensemble de quatre éléments en est le \nieme{3} ; pour finir – \italique{last but not least} –, faire préciser à Aristophane quel jour tombe ἡ μέση serait un détail inepte tant cela était connu de tous les Athéniens). Le texte avait été pourtant brillamment rétabli par August Nauck\nc{in \italique{Mélanges gréco-romains}, St.-Pétersburg (1866), t. II, pp. 468-71 ; voir aussi pp. 344-5 pour des témoignages sur les \italique{Thesmophories}.} et intégré avec explications dès l'édition aux Belles Lettres de Victor Coulon ; cela n'empêche pas la faute d'apparaître dans nombre d'éditions modernes. Venons-en donc à la bonne correction : l'apparition du τρίτη, \italique{troisième}, est due à la mauvaise lecture d'un Π' (abréviation de περί ou περ) en Γ' (\italique{troisième} lettre de l'alphabet) dans un ancien papyrus\nc{La corruption doit être ancienne puisqu'elle – conjointe au bavardage du scholiaste – rendait déjà perplexe Hésychius (τ 1440): τρίτη Θεσμοφορίων· ζητεῖται πῶς ἅμα μὲν λέγει τρίτην Θεσμοφορίων εἶναι, ἅμα δὲ μέσην, τεττάρων οὐσῶν ἡμερῶν.} – cette abréviation est fréquente, et source d'erreurs aussi dans l'autre sens ; ainsi a-t-il fallu en maintes occasions rétablir un `τρίτη' à partir d'un gamma\nc{Pour rester avec notre cher Aristophane, un `τρεῖς' a dû être restitué au vers 1507 des \italique{Guêpes} là où tous les manuscrits comportent un `γε'.}. Bref, voici l'histoire de notre début de \nieme{80} vers, très ingénieusement reconstituée : $$\overset{\text{\footnotesize{\allemand{Urtext}}}}{\text{ἐπείπερ ἐστὶ}}\xrightarrow{\text{onciales}}\text{ΕΠΕΙΠ'ΕCΤΙ}\xrightarrow{\text{corruption}}\text{ΕΠΕΙΓ'ΕCΤΙ}\xrightarrow{\text{interprétation}}\overset{\text{\footnotesize{\italique{Ravennas}}}}{\text{ἐπεὶ τρίτη ᾽στὶ.}}$$}
\cc{86–7}{νὴ τὸν Ποσειδῶ καὶ δί᾽ και᾽ ἂν πάθοις / ἀτὰρ τίν᾽ ἐκ ταύτης σὺ μηχανὴν ἔχεις;}{Ne voyant pas à quel nom féminin renvoyait `ἐκ ταύτης', on a retouché le texte du vers 87 à l'envi (ἐκ τούτων Bergler, εἰς ταύτας Sandbach, \&c.). 
C'est à mon avis aller un peu vite en besogne alors qu'il y a justement un problème textuel au vers précédent : il manque une syllabe, probablement au niveau du fautif `δί᾽ και᾽'. Je fais l'hypothèse que c'est justement à cet endroit que se trouvait ce nom féminin, repris par le `ἐκ ταύτης' du vers suivant. De quoi Euripide peut-il se mettre \italique{hors} en déployant quelque machine ? Réponse : tel le Strepsiade des \italique{Nuées}, c'est à la justice qu'il veut échapper. Au vu du texte conservé, la correction est évidente : il faut substituer au `δί᾽ και᾽' un `δικαίαν' ; la chute de l'`αν' s'expliquant par la courante erreur de «~saut du même au même~» (`δικαίαν' étant directement suivi de la particule `ἂν'). 
Lisons donc :
\citComm{νὴ τὸν Ποσειδῶ καὶ δικαίαν ἂν πάθοις.\\
ἀτὰρ τίν᾽ ἐκ ταύτης σὺ μηχανὴν ἔχεις;\\
\francais{Sainte Mer ! ce serait en justice que tu pâtirais.\\
Mais pour y échapper, de quelle machine disposes-tu ?}}
Cette correction est d'autant plus intéressante que ἡ δικαία, mot rare, est un synonyme poétique de δίκη\nc{Le mot est attesté dans l'\EtymMagnum à l'entrée Ἀθηνᾶ : le glosateur y recense Ἀθηναία, qui est à Ἀθήνη ce que δικαία est à δίκη (Ἀθήνη Ἀθηναία, ὡς δίκη δικαία).} ; Mnésiloque retourne-t-il contre Euripide une de ses propres sentences\nc{Tout au long de la pièce, Euripide est pris au piège de ses propres γνῶμαι (aux vers 194, 275-276, etc). Est-ce aussi le cas ici ? Quoi qu'il en soit, on trouve des expressions très proches aux vers 1379 des \italique{Nuées} : \citComm{Φειδιππίδης -- νὴ τὸν Δί᾽ ἐν δίκῃ γ᾽ ἄν,}ou 403 de \italique{Lysistrata} :\citComm{Πρόβουλος -- νὴ τὸν Ποσειδῶ τὸν ἁλυκὸν δίκαιά [δικαία ?] γε.}} ?\\
Wilamowitz\nc{Au commentaire du vers 681 de \italique{La Folie d'Héraclès} : second volume de son \italique{Euripides Herakles erklärt von Ulrich von Wilamowitz-Moellendorff}. Wilamowitz cite \italique{en passant} le vers 86 des \italique{Thesmophories} ainsi corrigé sans même préciser la correction -- dont il semble bien être l'auteur -- tant elle devait lui sembler aller de soi.} avait déjà proposé cette émendation sans qu'elle soit jamais reprise dans les éditions publiées depuis -- à une notable exception : celle de Prato.}
\cc{90–1}{ἐκκλησιάσοντ᾽ ἐν ταῖς γυναιξὶν κἂν δέῃ / λέξονθ᾽ ὑπὲρ ἐμοῦ.}{Je reprends la correction de Markland : substituer χἂν ($=$ καὶ ἃ ἂν) à κἂν ($=$ καὶ ἐὰν) qui présente un sens bien meilleur : Euripide ne dit certes pas « qu'au besoin » ou « s'il le faut » (κἂν) Agathon interviendra mais bien plutôt qu'il lui faut intervenir autant qu'il le pourra (χἂν). Euripide a sa vie en jeu – si l'on peut s'exprimer ainsi – et sait bien le courroux de la gent féminine ; l'inaction aurait forcément une issue fatale... heureuse pour Athènes et les Athéniennes mais un peu plus fâcheuse pour sa personne.}
\cc{100}{μύρμηκος ἀτραπούς.}{Willems, en accord avec le scholiaste, se propose d'expliquer ainsi l'expression μύρμηκος ἀτραπούς : « par allusion à la marche saccadée et tortueuse des fourmis. Image familière aux Anciens. Dans l'\italique{Anthologie}, XI, 78, il est question des galeries obliques et droites, λόξα καὶ ὀρθά, d'une fourmilière. Voir aussi Plutarque, \italique{de Solertia animal.}, II, p. 968a. Plaute a dit : \latin{formicinum gradum movere}. » Taillardat (§784) abonde en son sens avec d'autres exemples à l'appui. Cependant, pour suggérer l'image de \italique{sentiers de fourmis}, n'aurait-il pas fallu, en français comme en grec, un pluriel (μυρμήκ\textbf{ων} ἀτραπούς) ?\\
Hésychius ne le comprenait pas autrement et voyait dans `μύρμηκος' un nom propre (μ 1904) :
\citComm{Μύρμηκος ἀτραπούς· Ἀθήνησιν ἐν Σκαμβωνιδῶν ἔστι Μύρμηκος ἀτραπός, ἀπὸ ἥρωος Μύρμηκος ὀνομαζομένη.\\
\francais{Dans [le dème] athénien des Scambônides [σκαμβός signifie \italique{courbe}, \italique{tortu}] se trouve un sentier de Myrmex, qui tire son nom du héros Myrmex.}}}
